\section*{Chapter \ref{ch:iv2:risk-validation-and-regulation} --- Risk, Validation and Regulation}
\markboth{\protect\omTocMark{}{Answers: Risk, Validation and Regulation}}{\protect\omTocMark{}{Answers: Risk, Validation and Regulation}}

\begin{solution}{iq:iv2:risk-validation-and-regulation:1}
It is about \$2.33 million for one day, and about \$7.36 million for ten days.

The dollar volatility is \$1 million, and the 1\% quantile is 2.326 standard deviations. For ten days, multiply by $\sqrt{10} = 3.162$.

\iqwrong{Saying ``we can't lose more than \$2.33 million.'' VaR is the point the loss passes on one day in a hundred. It says nothing about how far the loss goes on that day.}
\end{solution}

\begin{solution}{iq:iv2:risk-validation-and-regulation:2}
It is about \$2.10 million, and diversification saves about \$0.11 million.

The dollar volatilities are 0.5 and 0.45. Then
\[
\sigma = \sqrt{0.25 + 0.2025 + 2 \times 0.8 \times 0.5 \times 0.45} = 0.901.
\]
And $0.901 \times 2.326 = 2.10$. Adding the standalone numbers gives $0.95 \times 2.326 = 2.21$.

\iqwrong{Multiplying the correlation into the volatilities and forgetting the factor 2 on the cross term. Weight it by both dollar volatilities and count it twice.}
\end{solution}

\begin{solution}{iq:iv2:risk-validation-and-regulation:3}
Expected shortfall is the average loss given that the loss is beyond VaR. For a standard normal, it is 2.338 at 97.5\% and 2.665 at 99\%. So 97.5\% shortfall is almost the same as the 99\% VaR of 2.326.

The formula is $\mathrm{ES}_{\alpha} = \varphi(z)/(1 - \alpha)$, with $z$ the quantile. So $\varphi(1.96)/0.025 = 2.338$.

That is why the regulators' switch to 97.5\% shortfall holds capital roughly level for a normal book and asks for more only if the tails are fatter.

\iqwrong{Thinking shortfall is a much higher number at any level. At 97.5\%, on a normal, it only barely exceeds the old 99\% VaR.}
\end{solution}

\begin{solution}{iq:iv2:risk-validation-and-regulation:4}
The 99\% VaR is about 5, with shortfall about 6.6.

One percent of 250 is 2.5 observations. Taking the third-worst as VaR is the conservative convention, and it gives 5. Interpolating between the second and third gives 5.5.

The shortfall averages the worst 2.5 observations: $(8 + 6 + 0.5 \times 5)/2.5 = 6.6$.

\iqwrong{Treating a single number as the answer without naming the convention. With so few points, the convention alone moves VaR by 10\%.}
\end{solution}

\begin{solution}{iq:iv2:risk-validation-and-regulation:5}
It is in the yellow zone, and the multiplier goes from 3 to 3.65.

Green is 0 to 4 exceptions, yellow is 5 to 9 with add-ons of 0.40, 0.50, 0.65, 0.75 and 0.85, and red is 10 or more, with 1.00.

\iqwrong{Saying that with 2.5 expected, the model has failed. Yellow means ``could go either way'', and supervisors look at why the exceptions came.}
\end{solution}

\begin{solution}{iq:iv2:risk-validation-and-regulation:6}
The half-life is about 11.2 days, and the new estimate is about 1.086\%.

The half-life solves $0.94^k = 0.5$, so $k = \ln 0.5/\ln 0.94 = 11.2$.

The variance is $0.94 \times 0.0001 + 0.06 \times 0.0004 = 0.000118$, and its root is 1.086\%.

\iqwrong{Updating the volatility rather than the variance: $0.94 \times 1\% + 0.06 \times 2\% = 1.06\%$. That is a different number, and it is not the model.}
\end{solution}

\begin{solution}{iq:iv2:risk-validation-and-regulation:7}
Total volatility is about 26.9\%, $R^2$ is about 45\%, and the portfolio of 100 has about 18.1\%.

Systematic volatility is $1.2 \times 15\% = 18\%$. Total is $\sqrt{0.18^2 + 0.20^2} = 0.269$. The $R^2$ is $0.0324/0.0724 = 44.8\%$.

In the portfolio, idiosyncratic volatility falls to $20\%/\sqrt{100} = 2\%$. So the total is $\sqrt{0.18^2 + 0.02^2} = 18.1\%$. Beta stays.

\iqwrong{Thinking diversification lowers it towards zero. It removes the idiosyncratic part only. The market risk does not diversify.}
\end{solution}

\begin{solution}{iq:iv2:risk-validation-and-regulation:8}
It is 50\%, from 130 to 65.

Keep the running peak. At each point, compute the fall from the peak, and keep the largest. That is one pass with two variables.

\iqwrong{Using the highest and lowest points overall. The low has to come after the peak, and here that happens to be the same pair, but a low before the high does not count.}
\end{solution}

\begin{solution}{iq:iv2:risk-validation-and-regulation:9}
It is about \$2.33 per option. That is too high, as gamma helps a long option.

The move is $2.326 \times \$2 = \$4.65$, and half of that is the delta loss. With gamma, $-0.5 \times 4.65 + 0.5 \times 0.02 \times 4.65^2 = -2.11$.

For a short option, the error runs the other way, which is worse: delta-normal then understates the loss.

\iqwrong{Dropping gamma on the view that it is small. On a move of 2.3 standard deviations it is 10\% of the number here, and for short options near expiry it can be most of the risk.}
\end{solution}

\begin{solution}{iq:iv2:risk-validation-and-regulation:10}
It is about \$698,000.

$\$50{,}000 \times 6 \times 2.326 = \$697{,}800$.

\iqwrong{Mixing up a yield's volatility in basis points with its volatility as a percentage of the level. A yield volatility quoted as a percentage needs the level to convert to basis points.}
\end{solution}

\begin{solution}{iq:iv2:risk-validation-and-regulation:11}
Each bond's 95\% VaR is 0. The pair's is 100. Shortfall at 95\% is 80 each and 103.2 for the pair. So VaR is not subadditive: combining two bonds raised it.

Each bond's loss has probability 4\%, which is within the 5\% tail, so its 5\% quantile is still zero. For the pair, the chance of at least one default is $1 - 0.96^2 = 7.84\%$, which is above 5\%, so VaR is 100.

For the shortfall, each bond's tail of 5\% holds 4\% at 100 and 1\% at 0, so $0.04 \times 100/0.05 = 80$. The pair's tail holds 0.16\% at 200 and 4.84\% at 100, so it is $(0.32 + 4.84)/0.05 = 103.2$, below $160$.

\iqwrong{Treating subadditivity as a mere technical concern. Here a VaR limit rewards concentration: you pick up rare-loss risks that sit just beyond the quantile, and the number stays at 0.}
\end{solution}

\begin{solution}{iq:iv2:risk-validation-and-regulation:12}
Yes, you reject: the likelihood ratio is 5.50, above the 3.84 cut-off, with a $p$-value of 1.9\%. The exact tail is 1.37\%. At 4 exceptions, the ratio is 0.77 and you do not reject. At 0, it is 5.03, and you reject the other way.

\[
\mathrm{LR} = -2\ln\frac{(1 - p)^{n - x}p^{x}}{(1 - \hat p)^{n - x}\hat p^{x}},
\]
with $\hat p = x/n$, compared with a $\chi^2_1$. The exact figure, $P(X \ge 7)$ for a binomial with $n = 250$ and $p = 1\%$, is a check on the asymptotic test.

The zero case surprises people. With no exceptions in 250 days, $\mathrm{LR} = -2 \times 250\ln(0.99) = 5.03$. The VaR is too large. That costs capital, even if it is a safer error.

\iqwrong{Using only one side. The test is two-sided, and an overly cautious model has failed too. Yet the traffic light punishes only the high side.}
\end{solution}

\begin{solution}{iq:iv2:risk-validation-and-regulation:13}
The transition probabilities are 2.06\% after a quiet day and 42.9\% after an exception, the likelihood ratio is 12.5, and independence is firmly rejected, with $p \approx 0.0004$.

$\pi_{01} = 5/243$, $\pi_{11} = 3/7$, and the pooled rate is $8/250 = 3.2\%$. The ratio is
\[
-2\bigl[\ln L(\text{pooled}) - \ln L(\pi_{01}, \pi_{11})\bigr] = 12.5,
\]
against a $\chi^2_1$ cut-off of 3.84.

Exceptions in clusters mean that VaR does not react to changes in volatility. So when a bad day comes, more follow, right when you need the number to be right.

\iqwrong{Passing the model because it looks fine in total. The count test does not care when they happen, so a model can pass on count with every exception in one week.}
\iqfollow{What does the conditional coverage test do?}{It adds the Kupiec and independence statistics, and compares the sum with a $\chi^2_2$.}
\end{solution}

\begin{solution}{iq:iv2:risk-validation-and-regulation:14}
By about 20\%: the ten-day volatility is 3.79 daily volatilities, not 3.16.

For an AR(1),
\[
\mathrm{Var}\Bigl(\sum_{i=1}^{m} r_i\Bigr) = \sigma^2\Bigl(m + 2\sum_{k=1}^{m-1}(m - k)\rho^k\Bigr).
\]
At $m = 10$ and $\rho = 0.2$, the bracket is 14.375, and the root is 3.79. The ratio is $3.79/3.16 = 1.20$.

At a negative autocorrelation, the rule overstates: mean-reverting returns accumulate less. And the formula gives $\sqrt{m}$ again at $\rho = 0$.

Also, volatility itself is not constant, and the rule ignores that too.

\iqwrong{Applying $\sqrt{10}$ to everything. Illiquid assets with smoothed prices tend to show positive autocorrelation, and for those it is exactly the wrong direction.}
\end{solution}

\begin{solution}{iq:iv2:risk-validation-and-regulation:15}
VaR is about 2.65, and shortfall about 3.69, against 2.33 and 2.67 for the normal.

A $t_\nu$ has variance $\nu/(\nu - 2) = 2$. So the unit-variance scale is $\sqrt{1/2}$, and $3.747 \times 0.7071 = 2.65$.

For shortfall, the closed form is
\[
\mathrm{ES} = \frac{f_\nu(q)}{0.01}\cdot\frac{\nu + q^2}{\nu - 1}\cdot\sqrt{\frac{\nu - 2}{\nu}} = 3.69.
\]

So at the same variance, VaR is 14\% higher, and shortfall 39\% higher. Fat tails hurt shortfall more.

\iqwrong{Forgetting to scale to unit variance. With the raw 3.747, you would compare a distribution of variance 2 with one of variance 1.}
\end{solution}

\begin{solution}{iq:iv2:risk-validation-and-regulation:16}
It is about 0.037 standard deviations, or about 1.6\% of VaR. For 0.5\%, you need about 103,000 paths.

The asymptotic standard error of a quantile is $\sqrt{p(1 - p)/n}/f(q)$. Here $\sqrt{0.0099/10{,}000} = 0.000995$, and $f(2.326) = 0.0267$. So $0.0373$, which is $0.0373/2.326 = 1.6\%$.

Since it falls as $\sqrt{n}$, going from 1.6\% to 0.5\% needs $(1.6/0.5)^2 \approx 10.3$ times more paths.

And a shortfall has a different standard error, and fatter tails need more paths for the same precision.

\iqwrong{Using the standard error of a mean, $1/\sqrt{n}$. The quantile estimator depends on the density at that point, and the density is low out in the tail.}
\end{solution}

\begin{solution}{iq:iv2:risk-validation-and-regulation:17}
The contributions are 1.110 and 0.987, adding up to 2.097.

The covariance of each with the total, divided by the total's volatility, gives $(0.25 + 0.18)/0.901 = 0.477$ and $(0.2025 + 0.18)/0.901 = 0.424$. Multiply by 2.326.

They add up by \index{Euler allocation}Euler's theorem. Volatility and normal VaR are homogeneous of degree 1 in position size, so the sum of size times marginal risk equals the total.

To confirm: bump the first position by 1\%, recompute VaR, and the change should be 1\% of its contribution, 0.0111.

\iqwrong{Allocating by standalone VaR. Those sum to 2.21, not to the total, and they give no sign of who adds risk at the margin. A hedge would get a positive allocation.}
\end{solution}

\begin{solution}{iq:iv2:risk-validation-and-regulation:18}
The loss is \$10.5 million, and equities need to fall a further 9\% for a loss of \$15 million.

Equities: $50 \times -20\% = -\$10$ million. Rates: $+\$40{,}000 \times 100 = +\$4$ million. Spreads: $-\$30{,}000 \times 150 = -\$4.5$ million.

The remaining \$4.5 million on \$50 million of equities is 9\%, so $-29\%$.

\iqwrong{Assuming rates always help. Whether a rate fall comes with an equity fall is exactly the assumption being stressed. In an inflation scare, both legs lose together, and the stress should have that scenario too.}
\end{solution}

\begin{solution}{iq:iv2:risk-validation-and-regulation:19}
The volatility is about 23.9\%, and factors explain about 81\% of the variance.

Portfolio factor exposures are $b = B^\top w = (1.04, 0.19)$. The factor variance is
\[
1.04^2 \times 0.04 + 2 \times 1.04 \times 0.19 \times 0.006 + 0.19^2 \times 0.01 = 0.0460.
\]
Idiosyncratic variance is $0.25 \times 0.03 + 0.09 \times 0.02 + 0.04 \times 0.04 = 0.0109$. The total is 0.0569.

\iqwrong{Using weights in place of squared weights for idiosyncratic variance. Those parts are independent, so they add with squared weights, which is why diversification works.}
\end{solution}

\begin{solution}{iq:iv2:risk-validation-and-regulation:20}
It is about \$208,000 at 0.9, and about \$465,000 at 0.5. So a correlation change raises VaR by 2.2 times with no change in positions.

The spread's volatility is $\$200{,}000\sqrt{2 - 2\rho}$, and multiply by 2.326.

So a low-VaR pair is long correlation: it relies on the correlation staying high. That correlation is lowest when you need it most: a crisis can pull a pair apart, or a takeover can break it in an afternoon. Recent correlation makes a position look safest just before it blows up, and so sizing by that number builds a large position at the wrong time.

\iqwrong{Treating gross of \$20 million and a VaR of \$208,000 as a small risk. Look at gross and stress the correlation, as gross is what you have to get out of.}
\end{solution}

\begin{solution}{iq:iv2:risk-validation-and-regulation:21}
Both have about the same VaR: 2.33 against 2.57. Their shortfalls differ by a factor of ten: 2.67 against 26.4.

For B, the loss of 50 is half of the 1\% tail. So VaR is the point where the normal part supplies the other 0.5\%: $0.995\,\Phi(-q) = 0.005$, so $q = 2.57$. The shortfall is
\[
\frac{0.005 \times 50 + 0.995\,\varphi(q)}{0.01} = 26.4.
\]

So shortfall sees the size of the tail, and is subadditive, which is why regulators moved to it.

\textbf{What you lose.} Shortfall is harder to backtest, as a single number doesn't pin down what you are checking. And it depends on few data points in the deep tail, so its estimate is less stable. It isn't elicitable alone: no scoring function has it as its best forecast, though it is jointly with VaR.

\iqwrong{Saying shortfall solves tail risk. It still relies on the same data, and if the 0.5\% event isn't in your sample, the number is 2.67.}
\end{solution}

\begin{solution}{iq:iv2:risk-validation-and-regulation:22}
It is about \$710,000 with five equal parts, and about \$485,000 if you can sell in a day.

One-day VaR is $\$10 \text{ million} \times 2\% \times 2.326 = \$465{,}000$.

Selling in equal parts, the position held across day $k$ is $(6 - k)/5$. So the total variance is $\sum_{k=1}^{5}((6 - k)/5)^2 = 2.2$ daily variances, a factor of $\sqrt{2.2} = 1.483$. That gives \$690,000. The general factor is $\sqrt{(n + 1)(2n + 1)/(6n)}$.

The half-spread of 20 basis points on the \$10 million adds \$20,000.

\iqwrong{Using $\sqrt{5}$. That holds the entire position for five days, but you are selling on the way, so the factor is 1.48 and not 2.24. And spread is a sure cost, so it adds in directly.}
\iqfollow{What is missing?}{Market impact, which grows with the trading rate. Selling quickly reduces this risk but costs more impact, and that trade-off is the optimal-execution problem.}
\end{solution}

\begin{solution}{iq:iv2:risk-validation-and-regulation:23}
Start by defining what it is for, then check the theory, rebuild it independently, test its limits, compare with outcomes, and approve only within a stated range of use. I would refuse if I could not reproduce its numbers or if it is being used outside its assumptions.

\begin{enumerate}
\item \textbf{Scope.} What it prices, its users, how much P\&L and capital depend on it. That sets the depth of review.
\item \textbf{Theory.} Are its assumptions right for the product? A one-factor model on a product paying on the correlation of two rates fails here, whatever the code does.
\item \textbf{Independent reimplementation.} Build it from the documentation and compare. This finds what the developer can't see.
\item \textbf{Benchmarks and limits.} Closed-form cases, parity, zero and huge volatility, and very short maturity. Check convergence.
\item \textbf{Inputs.} Where market data comes from, how it is calibrated, how stable the calibration is day to day.
\item \textbf{Outcomes.} P\&L explain, hedging simulation, and how far the model's price is from traded prices and consensus services.
\item \textbf{Conditions.} Limits on use, reserves for the uncertainty, a review date and monitoring.
\end{enumerate}

\begin{dated}{2026-11}{The guidance}\label{dat:iv2:risk-validation-and-regulation:mrm}
In the US, SR 26-2 of April 2026 superseded the earlier model-risk guidance, SR 11-7, and it keeps effective challenge by objective experts, with outcomes analysis and benchmarking in validation. In the UK, the PRA's SS1/23 has been in effect since May 2024, with five principles. Check the current text when you cite one.
\end{dated}

\iqwrong{Treating validation as a re-run of the developer's tests. Those show the code does what the developer meant. Validation asks whether that was the right thing to mean.}
\iqfollow{The desk says the trade must go today, and validation takes weeks. What do you do?}{Offer a conditional approval with a position limit, a reserve, and a date, and write it down. The worst outcome is a model that goes live with no one saying what it may do.}
\end{solution}

\begin{solution}{iq:iv2:risk-validation-and-regulation:24}
If the model is right, it stays green with probability 89.2\%. At 2\%, it stays green 43.9\% of the time, and goes red 3.0\%. At 3\%, green 12.8\%, and red 22.1\%.

Green is at most 4 exceptions, and red is at least 10. Binomial with $n = 250$. So $P(X \le 4 \mid p = 1\%) = 0.892$. At $p = 2\%$, $P(X \le 4) = 0.439$, and $P(X \ge 10) = 0.030$.

So a model that is wrong by a factor of two passes green nearly half the time. And a right model goes yellow 11\% of the time. That is why a one-year count is a weak test. To separate 1\% from 2\% cleanly, you need much more data or a test using more than counts.

\iqwrong{Treating green as a proof the model is right. A green zone is compatible with a doubled exception rate.}
\end{solution}

\begin{solution}{iq:iv2:risk-validation-and-regulation:25}
The capital is 109.5. Yesterday's VaR would have to be above 36.5 to matter. Stressed VaR adds 73, twice the 36.5 from the rest.

$\max(12, 3.65 \times 10) = 36.5$. $\max(18, 3.65 \times 20) = 73$. So the sum is 109.5.

Since the multiplier is above 3, the average nearly always wins. For the first term, a jump to more than 3.65 times the recent average would be needed.

Stressed VaR uses the current portfolio on inputs from a period of financial stress over 12 months. It doesn't fall in calm periods, and so capital is less procyclical than with VaR alone.

\begin{dated}{2026-11}{Where the rules are}\label{dat:iv2:risk-validation-and-regulation:frtb}
The January 2019 market-risk standard replaces this with 97.5\% expected shortfall. Adoption varies by place and date. As of mid-2026, the European Commission had adopted time-limited amendments applying from 1 January 2027, and the UK's PRA set internal-model changes for 1 January 2028. Check the current position before you cite it.
\end{dated}

\iqwrong{Thinking capital tracks yesterday's VaR. A sudden one-day jump has almost no effect, so it is the average that a desk must manage.}
\end{solution}

\begin{solution}{iq:iv2:risk-validation-and-regulation:26}
Delta-normal VaR is zero. The exact 99\% VaR is about \$6,635. The mean-and-variance normal gives about \$4,290, about 35\% too low.

The loss is $0.5 \times 500 \times x^2$ with $x = 2z$, so $1{,}000 z^2$. Its 99\% point comes from $z^2$ as a $\chi^2_1$, whose 99\% quantile is $6.635$. So the VaR is \$6,635. Both tails of the stock count, and $z = 2.576$ is the two-sided 1\% point.

The approximation has mean 1,000 and standard deviation $1{,}000\sqrt{2} = 1{,}414$. So $1{,}000 + 2.326 \times 1{,}414 = 4{,}290$.

The P\&L is strongly skewed, so a normal shape misses the tail.

\iqwrong{Reporting no risk since delta is zero. This book has no first-order risk at all, and its risk is all in the second order, which is what the method ignores.}
\end{solution}

\begin{solution}{iq:iv2:risk-validation-and-regulation:27}
The Cornish-Fisher quantile is about $-3.30$. On the $t_6$, the formula gives about $-3.03$, where the truth is $-2.57$. So it has overshot by 18\%. It is a series, and not a safe answer in fat tails.

\[
z_{\mathrm{CF}} = z + \frac{(z^2 - 1)S}{6} + \frac{(z^3 - 3z)K}{24} - \frac{(2z^3 - 5z)S^2}{36}.
\]
With $z = -2.326$, $S = -0.5$ and $K = 3$, that gives $-3.30$.

The $t_6$ has variance $6/4$, so its unit-variance quantile is $-3.143 \times 0.8165 = -2.57$. Its excess kurtosis is 3, and $S = 0$, so the formula gives $-2.326 - 0.70 = -3.03$.

So the same kurtosis can mean different tails. The expansion is exact only for small departures from normal, and it can even produce non-monotone quantiles at higher skew and kurtosis. It is a useful correction for small $S$ and $K$. With fat tails, fit a distribution or use the empirical tail.

\iqwrong{Assuming a higher-moment correction is always closer to the truth. Here it is further from the truth than the plain normal in size: 0.46 too far against 0.24 too near. And sample skew and kurtosis are themselves noisy.}
\end{solution}

\begin{solution}{iq:iv2:risk-validation-and-regulation:28}
The mean is 0 if the model is right, and about $-1.16$ if the true volatility is 1.2 times the model's. A 5\% cut-off is about $-0.70$, by simulation, and that detects the 1.2 case about 77\% of the time over 250 days.

Under the model, each day's expectation is $E[X I] = -\alpha\,\mathrm{ES}$, so $Z$ has mean $1 - 1 = 0$.

If the true volatility is $s$, $E[X I] = -s\,\varphi(v/s)$, with $v = 1.96$. So the mean of $Z$ is $1 - s\,\varphi(v/s)/(0.025\,\mathrm{ES})$, and at $s = 1.2$ that gives $-1.158$.

The statistic has one nice property for a bank: since it uses sizes as well as counts, a too-small shortfall shows up even if VaR is right. The cut-off has to come from simulation, and it depends on the assumed model, though it is stable for a range of tails.

\iqwrong{Thinking you can't backtest shortfall at all. Its lack of elicitability rules out one type of comparison, not all tests. What it needs is the model's own forecast, day by day.}
\end{solution}

\begin{solution}{iq:iv2:risk-validation-and-regulation:29}
It is about 12\% a day at first, so about 2.2 exceptions over 20 days, against 0.2 expected. The exponentially weighted model is at about 3.6\% on the tenth day.

At first, the threshold is still 2.326 old volatilities, which is 1.163 new ones, and $\Phi(-1.163) = 12.2\%$. Held fixed for 20 days, that would give 2.45 exceptions. But the new days enter the window as the worst ones, and each lifts the third-worst a little. In a simulation, the rate falls from about 12.7\% to about 9\% over the 20 days, and the total is about 2.2.

For the exponential model, after 10 days a share $1 - 0.94^{10} = 46\%$ of the weight is on the new regime. So variance is $0.54 + 0.46 \times 4 = 2.38$ old units, and volatility is 1.54 old units. The threshold is $2.326 \times 1.544/2 = 1.80$ new volatilities, so $\Phi(-1.80) = 3.6\%$.

So historical VaR's exceptions cluster, and they come right after volatility jumps. An exponential model catches up, though it still starts behind. And the historical model's VaR keeps the high number for a year after the storm ends.

\iqwrong{Raising the confidence level to solve it. A higher quantile is a higher level, but it adapts no faster.}
\end{solution}

\begin{solution}{iq:iv2:risk-validation-and-regulation:30}
The VaR is about 3.90, and the shortfall about 5.62. As the shape reaches 1, shortfall becomes infinite.

With $u = 2$, $\beta = 1$, $\xi = 0.2$, and exceedance rate 5\%,
\[
\mathrm{VaR}_p = u + \frac{\beta}{\xi}\Bigl(\Bigl(\frac{1 - p}{0.05}\Bigr)^{-\xi} - 1\Bigr) = 2 + 5(0.2^{-0.2} - 1) = 3.90.
\]
And
\[
\mathrm{ES}_p = \frac{\mathrm{VaR}_p + \beta - \xi u}{1 - \xi} = \frac{3.90 + 1 - 0.4}{0.8} = 5.62.
\]

The $1 - \xi$ in the denominator is why. A shape of 1 or more means the tail has no finite mean, so no average loss exists to report. At a shape of 0, the tail is exponential: VaR is $2 + \ln 5 = 3.61$, and shortfall exceeds it by exactly $\beta$.

\iqwrong{Choosing a threshold so high that few points remain, or so low that the Pareto shape doesn't hold. Plot the estimate of $\xi$ against the threshold, and take a stable stretch.}
\end{solution}

\begin{solution}{iq:iv2:risk-validation-and-regulation:31}
It is about 0.13\% under the Gaussian copula, and about 0.29\% under the $t_4$, against 0.01\% for independence. The limit of tail dependence is 0 for the Gaussian and about 0.25 for the $t_4$.

The Gaussian figure is the bivariate normal at $(c, c)$, with $c = -2.326$, and the $t$ figure is by simulation. So given that one falls past its 1\% point, the other follows with probability about 13\% or 29\%.

The tail dependence coefficient for a $t$ copula is
\[
\lambda = 2\,t_{\nu + 1}\Bigl(-\sqrt{\frac{(\nu + 1)(1 - \rho)}{1 + \rho}}\Bigr) = 0.253.
\]
For the Gaussian, $\lambda = 0$ whenever $|\rho| < 1$: in the limit, the extremes are independent.

So the same correlation can mean twice the joint crash risk, and a Gaussian copula with fat marginals still misses it.

\iqwrong{Thinking fat-tailed marginals fix joint tails. They change each asset's own tail, but it is the copula that decides whether they crash together.}
\end{solution}

\begin{solution}{iq:iv2:risk-validation-and-regulation:32}
The capital measure is 16. It is not $10\sqrt{12} = 34.6$, because each horizon's scaling applies only to the factors with that horizon, and factors are shocked over the extra time alone.

\[
\mathrm{ES} = \sqrt{\mathrm{ES}_T(P)^2 + \sum_{j \ge 2}\Bigl(\mathrm{ES}_T(P, j)\sqrt{\frac{\mathrm{LH}_j - \mathrm{LH}_{j-1}}{T}}\Bigr)^{2}}.
\]
The terms are $10^2$, $(8\sqrt{1})^2 = 64$, $(5\sqrt{2})^2 = 50$, $(3\sqrt{2})^2 = 18$ and $(2\sqrt{6})^2 = 24$. The total is 256, so 16.

If only 10-day factors exist, the later terms vanish, and it is just 10. If every factor is 120-day, all are 10, and the sum is $100(1 + 1 + 2 + 2 + 6) = 1{,}200$, so $10\sqrt{12} = 34.6$. So that naive number is the upper limit, and the formula passes this check.

\iqwrong{Scaling the whole portfolio to its longest horizon. That is the top limit, and it charges liquid factors for an illiquid horizon.}
\end{solution}

\begin{solution}{iq:iv2:risk-validation-and-regulation:33}
Volatility and VaR are halved. You would then expect about 31 exceptions a year in a 99\% test, against 2.5. A validator catches it by reperforming one small calculation by hand before any data run.

For small changes, $x_t + x_{t-1} \approx 2x_{t-1}$, so the ratio is half the intended one, and every return is halved. So every volatility is halved, and VaR with it. It is exact to first order, whatever the volatility.

For normal returns, the threshold is $2.326/2 = 1.163$ true standard deviations. So the rate is $\Phi(-1.163) = 12.2\%$, or 31 days in 250.

\textbf{How to catch it early.} Hand-compute one number. Compare the input volatility with an independent source. Put in a series with a known volatility and check that it comes back. Don't let the production model live in an unreviewed spreadsheet with manual steps.

\iqwrong{Trusting the backtest to catch it. If a change of model comes with a quiet market, or hedged positions, the exceptions can take months to show. And a model that lowers VaR is rarely questioned by the people it helps.}
\end{solution}

\begin{solution}{iq:iv2:risk-validation-and-regulation:34}
It is 21 at $V - C = 5$, and about 5.34 at $-20$. The floor on the multiplier is 5\%, which gives 0.7 when RC is zero.

At $V - C = 5$, $\mathrm{RC} = 5$, and the multiplier caps at 1: $1.4 \times (5 + 10) = 21$.

At $-20$, $\mathrm{RC} = 0$, and the multiplier is $0.05 + 0.95\,e^{-20/19} = 0.382$. So $1.4 \times 3.82 = 5.34$.

Being out of the money or over-collateralised reduces future exposure as well, as it would take a move to get back to positive. As $V - C \to -\infty$, the multiplier is $0.05$, and $1.4 \times 0.5 = 0.7$. The floor exists so that excess collateral never takes exposure to zero.

\iqwrong{Treating RC as the whole exposure. Exposure is also the potential change over the rest of the trade, and that's what the add-on stands for.}
\iqfollow{Why the 1.4?}{It is the supervisory scaling factor, $\alpha$, a fixed number in the rule, and it covers things like wrong-way risk and model error.}
\end{solution}

\begin{solution}{iq:iv2:risk-validation-and-regulation:35}
The most likely such scenario is $x_1 = -5.38$, $x_2 = -0.77$, and it is about 2.77 standard deviations from the origin.

Under normal factors, maximise the density given the loss, which means minimising $x^\top\Sigma^{-1}x$ given $b^\top x = -10$. With a Lagrangian,
\[
x^* = -\frac{10\,\Sigma b}{b^\top\Sigma b},
\]
with $b = (2, -1)$, $\Sigma b = (7, 1)$ and $b^\top\Sigma b = 13$. So $x^* = -(10/13)(7, 1)$. The distance is $10/\sqrt{13} = 2.77$. The P\&L has a standard deviation of 3.61, so 10 is the same 2.77 of them.

\textbf{The second factor.} You are short it, so you'd gain from a fall, and a fall in $x_2$ gives you $+0.77$. It falls anyway, since it is positively correlated with the first factor, and a large drop in that one drags it along. The scenario is what the correlation says is likely, and not what hurts most per unit.

Compare a stress on the first factor alone, $(-5, 0)$, which also loses 10. Its distance is 2.89, so it is less likely than $x^*$.

\iqwrong{Stressing one factor at a time and calling the worst a reverse stress. The point is to find the plausible combination, and it is often not the one you'd have guessed.}
\end{solution}

\begin{solution}{iq:iv2:risk-validation-and-regulation:36}
That matrix isn't a correlation matrix. Its smallest eigenvalue is $-0.547$. With the two pairs at 0.9, the (1, 3) correlation must lie in $[0.62, 1]$.

The range comes from a determinant condition, or from partial correlation:
\[
\rho_{13} \in \rho_{12}\rho_{23} \pm \sqrt{(1 - \rho_{12}^2)(1 - \rho_{23}^2)} = 0.81 \pm 0.19.
\]

If you use it anyway, there's a portfolio with negative variance: weight along the eigenvector with the negative eigenvalue, and the formula takes a square root of a negative number, or returns a risk of zero, so an optimiser will pile into it.

To repair it, change as little as possible: set negative eigenvalues to a small positive and rescale to a unit diagonal, or solve a nearest-correlation problem. Then re-read the result, since the repair may have moved the numbers you meant to stress.

\iqwrong{Thinking any number between $-1$ and 1 will do. Pairwise bounds aren't enough, as correlations have to be mutually consistent. Stress tests that set a few pairs by hand break this all the time.}
\end{solution}

\begin{solution}{iq:iv2:risk-validation-and-regulation:37}
The mean stays near 2.5, but the variance is about 10.6 against 2.475. So it's about 4.3 times too large, and the count behaves as if you had only 58 independent observations.

The windows share nine of ten days, so an exception today tends to stay one for several days. The mean of a count is unaffected, as each window has a 1\% marginal chance. But exceptions come in clumps of up to about ten, which raises the variance.

\textbf{What it does.} Kupiec assumes a binomial. With the real, wider distribution, you reject a correct model far too often at the usual cut-off. And the independence test fails by construction.

\textbf{What to do.} Backtest at one day, which the rules do: one-day VaR for backtesting, ten-day for capital. Or use non-overlapping ten-day windows, though that leaves only 25 observations a year. Or simulate the null distribution of the overlapping count, and use its cut-off.

\iqwrong{Counting overlapping windows as 250 observations of ten-day risk. You have only 25 independent ones. Any statement about tail frequency from that has to be weak.}
\end{solution}

\begin{solution}{iq:iv2:risk-validation-and-regulation:38}
The Kupiec test passes, with a ratio of 2.61 and a $p$-value of 10.6\%. The independence test fails badly, with a ratio of 21.9 and a $p$-value of about 0.0003\%. So the count is acceptable and the timing is not. The model is too slow.

\textbf{Numbers.} With $\hat p = 9/500 = 1.8\%$, Kupiec gives 2.61 against a cut-off of 3.84, and the exact tail $P(X \ge 9) = 6.7\%$. The transitions give $\pi_{01} = 5/490 = 1.02\%$ and $\pi_{11} = 4/9 = 44.4\%$, a ratio of 21.9. An exception makes the next one more than 40 times likelier.

\textbf{What I would look at, in order.}
\begin{enumerate}
\item \textbf{Which P\&L?} Backtest on hypothetical P\&L with positions frozen, without fees and intraday trading. If exceptions come only on actual P\&L, the issue may be new trades and not the model.
\item \textbf{When?} Line up the clusters with market events. If they fit volatility jumps, the window is too slow: test a volatility-scaled historical simulation or an exponential weighting.
\item \textbf{Which positions?} Attribute each exception by factor. A single desk or a factor missing from the model (a basis, a vol skew, a new product) points to a coverage issue, not a slow model. The fix differs.
\item \textbf{Data.} Stale proxies and missing days smooth volatility and create breaches when the actual prices catch up.
\item \textbf{Mis-sized positions.} Check for gamma and jump risk near expiry, since a linear approximation can look right until a large move.
\item \textbf{Comparison.} Run the candidate model on the same history: the new one has to fix the clustering without just inflating the level.
\end{enumerate}

\textbf{What I would say.} To the desk: the level is fine, but the model is slow, and limits based on it are too loose just after a shock. To the regulator: exactly that, with the evidence, the remediation plan and date, and what we are doing in the meantime, such as a temporary add-on. On a 250-day clock, a cluster can also move the bank into the yellow zone, so raise it early rather than wait for the multiplier to change.

\iqwrong{Reporting the 2.61 and stopping there. On a count alone, this model passes. Risk models fail in the timing of the exceptions more often than in the total number.}
\iqfollow{Why not raise VaR by a flat factor until the count is right?}{That raises the level all the time, costing capital, and leaves the clustering in place. The exceptions still come together, just slightly less often.}
\end{solution}

\begin{solution}{iq:iv2:risk-validation-and-regulation:39}
Delta-normal says 0. Gamma alone gives a 99\% VaR of about \$127,000. Gamma and vega together give about \$174,000, with shortfall about \$223,000. Under Student-$t$ moves, the VaR is about \$252,000, and shortfall about \$525,000. The stress loses about \$1.18 million. I would set the limit on the stress, not the VaR.

\textbf{The P\&L.} With $x$ the stock's move and $\Delta v$ the volatility change,
\[
\mathrm{P\&L} = 22{,}500 - 10{,}000\,x^2 - 40{,}000\,\Delta v.
\]
The theta is fair: $\tfrac12 \times 20{,}000 \times 1.5^2 = 22{,}500$, so the expected P\&L is zero.

\textbf{Gamma alone.} $x^2 = 2.25 z^2$, whose 99\% point is $2.25 \times 6.635$. So the loss is $149{,}285 - 22{,}500 = \$126{,}800$.

\textbf{Vega alone.} $40{,}000 \times 2.326 = \$93{,}000$. The two added would be \$220,000. They are not perfectly linked, and vega hurts in only one tail, so the combined figure is \$174,000 by simulation of $\Delta v = -0.7 z + \sqrt{0.51}\,\varepsilon$. A falling stock brings rising vol, so down moves are hit twice and up moves partly offset.

\textbf{Fat tails.} Scale both shocks to a unit-variance $t_4$ with one shared scale. VaR rises by 45\%, and shortfall more than doubles. Gamma is a square, so it makes fat tails much worse.

\textbf{The stress.} $-\$10$ and $+5$ points: $10{,}000 \times 100 - 22{,}500 + 200{,}000 = \$1{,}177{,}500$. That is nearly seven times the normal VaR. A move of $-\$10$ is 6.7 standard deviations. On a normal, that never happens. In a single-stock book, an earnings surprise or takeover can do it in one jump.

\textbf{The limit.} Set it on scenarios: a grid of stock moves and vol moves, with a maximum loss in the grid, such as $\pm 10\%$ with $\pm 5$ to $10$ points. VaR doesn't see the gap, as it's a quadratic that grows with the square of the move. Then keep VaR as a daily measure, and add limits on gamma and vega by expiry.

\iqwrong{Reporting delta-normal for a delta-hedged book. The number that says ``no risk'' is the first one every risk manager should distrust. Quote what breaks it.}
\iqfollow{The stock jumps $+\$10$ and vol falls 3 points. What then?}{The loss is $1{,}000{,}000 - 22{,}500 - 120{,}000 = \$857{,}500$. An up gap still costs a lot, so the book is hurt both ways.}
\end{solution}

\begin{solution}{iq:iv2:risk-validation-and-regulation:40}
The model is correct at the money and in the normal range, is first-order accurate with an even-odd oscillation, is badly off in relative terms far out of the money, and has a real defect at low volatility. I would approve it with conditions once that defect is fixed.

\textbf{Oscillation.} The errors are $-0.0399$, $+0.0344$, $-0.0200$, $+0.0174$, $-0.0100$, $+0.0087$. The sign follows parity of $N$, as the final node either sits on the strike or straddles it. The size halves as $N$ doubles, so the order is $1/N$. \emph{Fix:} average two adjacent steps. For $N = 100, 101$ the average is 10.4493, an error of $-0.0013$, which beats $N = 1{,}000$ alone at $-0.0020$. Or smooth the last step by using the closed form there.

\textbf{Far out of the money.} At a strike of 180, $N = 50$ is $12.9\%$ low in relative terms and $N = 200$ is $3.6\%$ low. The absolute error is tiny, but a desk selling wings at a low price cares about the ratio. \emph{Fix:} a relative tolerance by strike, more steps there, or a closed form where it exists. State the range of use.

\textbf{Low volatility.} The up probability is $p = (e^{r\Delta t} - d)/(u - d)$, and it needs $\sigma > r\sqrt{\Delta t}$. At $\sigma = 0.1\%$, $N = 500$, $\Delta t = 0.002$: $r\sqrt{\Delta t} = 0.00224$, well above $\sigma$. So $p > 1$, and the weights are not probabilities. Backward induction then multiplies errors at every step until the result blows up. The limit should be $S - Ke^{-rT} = 4.88$. \emph{Fix:} check $0 \le p \le 1$ and refuse or re-grid, or use steps above 2,500 here. A trinomial or a drift-matched lattice avoids this.

\textbf{What else before signing.}
\begin{itemize}
\item Parity: call minus put against $S - Ke^{-rT}$ across the grid. At these inputs, it agrees to 12 digits.
\item Greeks: check bumped deltas and gammas for noise. A tree's gamma can oscillate even when the price looks fine.
\item Limits: zero rate, very short and very long expiry, huge volatility, dividends, deep in the money.
\item Exercise: if it is for American options, compare with a separate fine grid, as there's no closed form. Check that an American call with no dividends equals the European.
\item Discounting: a tree with one rate error is still close to the right answer. Test a case that singles out the discount, such as the no-arbitrage forward.
\item Speed and memory, and what happens on bad input.
\end{itemize}

\iqwrong{Comparing one price at one step count, seeing a match to two decimals and signing. A limit case is what catches this defect, since no typical input shows it. The worst failure is a large wrong number with no warning raised.}
\end{solution}
